% Источники: exam/materials/моделирование.pdf, раздел 6; exam/materials/преподаватель/2020/01-05-2020-Model_L11__1_.pdf; exam/materials/преподаватель/2020/01-05-2020-Model_L11__2_.pdf; exam/materials/прошлогодние ответы/Modelirovanie_1_-3.pdf, вопрос 36; GitHub HumsterProgrammer/iu9-notes/8sem/Моделирование/Моделирование_лекция-15_02-04-26.md.
% Примечание по расхождению: в exam/materials/моделирование.pdf E-оптимальность сформулирована через максимальную дисперсию предсказания, в Modelirovanie_1_-3.pdf — через максимальную дисперсию коэффициентов; ниже сохранена формулировка основного источника.
\section{Планирование эксперимента. Критерии оптимальности планов. Свойства планов.}

Критерии оптимальности планов и способы организации активного эксперимента можно разделить на 2 группы.

К I-й группе относят критерии, связанные с точностью оценок коэффициентов регрессии.

\begin{itemize}
  \item \textbf{D-оптимальность:} план считается D-оптимальным, если он минимизирует обобщенную дисперсию оценок параметров модели. Математическая формулировка: максимизация детерминанта информационной матрицы $(X^T X)$.
  \item \textbf{A-оптимальность:} план считается A-оптимальным, если он минимизирует сумму дисперсий оценок параметров модели. Математическая формулировка: минимизация следа обратной информационной матрицы $(X^T X)^{-1}$.
  \item \textbf{E-оптимальность:} план считается E-оптимальным, если он минимизирует максимальную дисперсию предсказания. Математическая формулировка: минимизация наибольшего собственного значения матрицы $(X^T X)^{-1}$.
\end{itemize}

Выбор критерия зависит от задачи исследования, так при изучении влияния отдельных факторов на поведение объекта применяют критерий Е-оптимальности, а при поиске оптимума функции отклика — D-оптимальности. Если построение D-оптимального плана вызывает затруднения, то можно перейти к А-оптимальному плану, построение которого осуществляется проще.

Ко II-ой группе относят критерии, которые определяют точность предсказания отклика с помощью построенной модели.

\begin{itemize}
  \item \textbf{G-оптимальность:} план считается G-оптимальным, если он минимизирует максимальную дисперсию предсказания по всем точкам пространства факторов.
  \item \textbf{Q-оптимальность / I:} план считается Q-оптимальным, если он минимизирует среднюю дисперсию предсказания по всем точкам пространства факторов.
\end{itemize}

Точка в факторного пространства, соответствующая нулевым уровням всех факторов, называется \textbf{центром плана}.

\textbf{Ортогональным} называется план, для которого выполняется условие парной ортогональности столбцов матрицы планирования. При ортогональном планировании коэффициенты полинома определяются независимо друг от друга — вычеркивание или добавление слагаемых в функции отклика не изменяет значения остальных коэффициентов полинома.

\textbf{Ротатабельность.} Использование ротатабельных планов обеспечивает для любого направления от центра эксперимента равнозначность точности оценки функции отклика (постоянство дисперсии предсказания) на равных расстояниях от центра эксперимента.

\textbf{Униморфность.} В ряде случаев при исследовании поверхности отклика требуется соблюдение постоянства значений дисперсии ошибки в некоторой области вокруг центра эксперимента. Выполнение такого требования целесообразно, когда исследователь не знает точно расположение области поверхности отклика с оптимальными значениями параметров.

По соотношению между числом неизвестных параметров и числом факторов все планы подразделяются на:
\begin{enumerate}
  \item ненасыщенные: число параметров $<$ числа факторов;
  \item насыщенные: число параметров $=$ числу факторов;
  \item сверхнасыщенные: число параметров $>$ числа факторов.
\end{enumerate}

Для некоторых планов важную роль играет свойство композиционности. Композиционные планы для построения полиномов второго порядка получают добавлением некоторых точек к планам формирования линейных функций.

\clearpage
